\section{Conclusion}
\label{sec:conclusion}

Open-ended Web exploration can support reusable functional model induction,
but discovery alone is not enough. An agent must distinguish a proposed
function from an executor-reported completion and from an outcome supported by
observable evidence. It must also account for the environmental consequences
of the interactions used to obtain that evidence. VERA connects these concerns
in a traceable model: it fixes functional claims before execution, records
context-conditioned action risk, preserves before--after evidence, and admits
only supported outcomes while retaining unresolved verification opportunities.

Across controlled studies, this design yields more precise knowledge admission,
extends evidence-supported coverage through persistent recovery, and improves
the contextual expression of risk consequences. The results are bounded: they
do not establish general exploration optimality or evaluated harm prevention.
They nevertheless support a broader view of trustworthy environment learning.
What an agent learns should be evaluated not only by how much functionality it
discovers, but also by what evidence makes those discoveries reusable and what
actions were taken to acquire that evidence.
